\section{Conditioning on survival}
\label{m:sec:qsd}

Absorption is often inevitable, as it is here for $p<\tfrac12$, and yet a process
on its way to absorption need not be featureless. In many cases it settles, long
before it is absorbed, into a distribution that no longer changes with time --- a
distribution describing the population conditional on its not yet having been
absorbed. The phenomenon is \emph{quasi-stationarity}, and the distribution is the
quasi-stationary, or Yaglom, limit \cite{Yaglom1947}; modern accounts are given by
Méléard and Villemonais \cite{Meleard2012} and by Collet, Martínez and San Martín
\cite{Collet2013}. This section establishes it for the binary Galton--Watson
process by computing the limiting conditional mean and variance, and in doing so
introduces the constant that \ChA{} is devoted to.

\begin{definition}[Quasi-stationary and Yaglom laws]
\label{m:def:qsd}
Let $\{X_n\}$ be a Markov chain with absorbing state $0$. A probability
distribution $\nu$ on the non-absorbing states is \emph{quasi-stationary} if
\begin{equation}
\label{m:eq:qsddef}
\pr{X_n = j\mid X_n\neq0} = \nu_j
\qquad(n\ge0,\ j\neq0)
\end{equation}
whenever $X_0\sim\nu$. If, from a specified initial state,
\begin{equation}
\label{m:eq:yaglomdef}
\pr{X_n = j\mid X_n\neq0}\longrightarrow\nu_j ,
\end{equation}
then $\nu$ is the corresponding \emph{Yaglom limit}.
\end{definition}

For a subcritical Galton--Watson process satisfying the usual non-degeneracy and
integrability conditions the Yaglom limit exists; the original result is Yaglom's
\cite{Yaglom1947}, with modern treatments in
\cite{Harris1963,Athreya1972,Meleard2012}. The binary offspring law satisfies those
conditions, and existence is taken from that literature rather than re-proved here.
What is computed below is the survival scale and the first two limiting moments,
explicitly and in terms of a single constant. The distinction matters: convergence
of two moments is consistent with the Yaglom theorem but does not on its own
establish convergence of the full conditional distribution.

\subsection{Late-time survival probability}
\label{m:sec:lateS}

Everything follows from the behaviour of $S_n$. Take $p<\tfrac12$, so that death of
an individual is more probable than its division. Then $S_n\to0$, and as it does so
the quadratic term in
\begin{equation*}
S_{n+1} = 2pS_n - pS_n^2
\end{equation*}
becomes negligible beside the linear one. The recursion is therefore increasingly
well approximated by
\begin{equation}
\label{m:eq:attenuated}
S_{n+1}\approx 2p\,S_n ,
\end{equation}
and so at late times
\begin{equation}
\label{m:eq:lateS}
S_n \sim A(p)\,(2p)^n\qquad(n\to\infty),
\qquad
A(p) = \lim_{n\to\infty}\frac{S_n}{(2p)^n} .
\end{equation}
The constant $A(p)$ absorbs the cumulative effect of the quadratic term over the
early generations, before the linear regime takes hold. It is precisely what
remains of the conditions imposed on $S_n$ before the recursion attenuates to
\cref{m:eq:attenuated}, which is why no argument confined to late times will produce
it. That the limit exists, is finite, and is strictly positive for
$p\in[0,\tfrac12)$ is proved in \ChA{} by way of an infinite product, which also
locates it in $(0,\tfrac12]$; the present section needs only the definition and the
existence.

\begin{remark}
\label{m:rem:asympnotidentity}
The relation \cref{m:eq:lateS} is an asymptotic equivalence, not an identity. The
limit of $S_n$ itself is $0$ for $p<\tfrac12$, and one must not write
$\lim_n S_n = A(p)(2p)^n$, whose right-hand side still depends on $n$. What is
asserted is that the ratio $S_n/(2p)^n$ converges.
\end{remark}

\subsection{The limiting conditional mean}
\label{m:sec:condmean}

Decompose the expected population size into the contributions from surviving and
from extinct realisations:
\begin{equation*}
\ex{Z_n} = \pr{Z_n=0}\ex{Z_n\mid Z_n=0} + \pr{Z_n>0}\ex{Z_n\mid Z_n>0} .
\end{equation*}
The first term vanishes identically, since $Z_n=0$ on that event, so with
$\pr{Z_n>0}=S_n$ and $\ex{Z_n}=(2p)^n$ from \cref{m:eq:meanZn},
\begin{equation}
\label{m:eq:condmeandef}
\ex{Z_n\mid Z_n>0} = \frac{\ex{Z_n}}{S_n} = \frac{(2p)^n}{S_n} .
\end{equation}
The numerator is known exactly; all the content is in $S_n$, and \cref{m:eq:lateS}
supplies it:
\begin{equation}
\label{m:eq:qsmean}
\ex{Z_n\mid Z_n>0} \;\xrightarrow[n\to\infty]{}\;
\frac{(2p)^n}{A(p)(2p)^n} = \frac{1}{A(p)} .
\end{equation}
The geometric factors cancel exactly. At late times the surviving population has a
constant expected size, and that size is the reciprocal of $A(p)$; accordingly
$1/A(p)$ is called the \emph{quasi-stationary mean} below.

\Cref{m:fig:condmean} shows the convergence, and shows at the same time the
difficulty that \ChA{} has to confront: the closer $p$ lies to $\tfrac12$,
the higher the plateau and the longer the process takes to reach it.

\begin{figure}[ht]
\centering
\includegraphics[width=0.78\textwidth]{conditionalMean.pdf}
\caption{The conditional mean $(2p)^n/S_n$ against generation number $n$, with
$S_n$ obtained by iterating \cref{m:eq:SnIter}. From bottom to top,
$p=0.3,\,0.4,\,0.45,\,0.48,\,0.49$. Each curve approaches the constant $1/A(p)$;
the limiting values are $2.825$, $4.208$, $6.893$, $14.703$ and $27.476$. The
lowest three have essentially converged by $n=100$, but the $p=0.49$ curve has
reached only $24.1$ of its limiting $27.5$ --- an early sign of a near-critical
slowdown that \ChA{} analyses and works around. The vertical axis is the
conditional mean.}
\label{m:fig:condmean}
\end{figure}

\subsection{The limiting conditional variance}
\label{m:sec:variance}

The mean alone does not establish quasi-stationarity. To show that the conditional
\emph{distribution} settles we need at least the second moment, and computing it is
a worthwhile exercise in its own right.

\subsubsection{The second moment}

For a \rv{} $X$ the moment generating function is $\ex{\ee^{tX}}$, which in terms
of the \pgf{} is $\phi(\ee^t)$. Write $M_n(t) = \phi_n(\ee^t)$ for the moment
generating function of $Z_n$, so that moments are recovered by
\begin{equation}
\label{m:eq:mgfmoments}
\ex{Z_n^{\,\varrho}} = \dda{^{\varrho}}{t^{\varrho}}M_n(t)\bigg|_{t=0} .
\end{equation}
For the binary offspring law \cref{m:eq:pgfbinary} the first-generation function is
\begin{equation}
\label{m:eq:mgf1}
M_1(t) = \phi(\ee^t) = p\,\ee^{2t} + (1-p),
\end{equation}
whose second derivative at $t=0$ is $4p$, so $\ex{Z_1^2}=4p$. Reaching generation
$n$ requires a little more. Since $\phi_n$ is the $n$-fold composition of $\phi$,
\begin{equation*}
M_n(t) = \phi\bigl(\phi(\dots\phi(\ee^t)\dots)\bigr),
\end{equation*}
and differentiating by the chain rule gives a product of $n+1$ factors,
\begin{equation}
\label{m:eq:dmgf1}
M_n'(t) = \phi'\bigl(\phi_{n-1}(\ee^t)\bigr)\cdot\phi'\bigl(\phi_{n-2}(\ee^t)\bigr)
\cdots\phi'\bigl(\phi_1(\ee^t)\bigr)\cdot\phi'(\ee^t)\cdot\ee^t .
\end{equation}
Now $\phi'(z) = 2pz$ for the binary law, so each factor collapses,
\begin{equation*}
\phi'\bigl(\phi_k(\ee^t)\bigr) = 2p\,\phi_k(\ee^t) = 2p\,M_k(t),
\qquad
\phi'(\ee^t) = 2p\,\ee^t ,
\end{equation*}
and \cref{m:eq:dmgf1} becomes
\begin{equation}
\label{m:eq:dmgf2}
M_n'(t) = (2p)^n\,M_{n-1}(t)\,M_{n-2}(t)\cdots M_1(t)\,\ee^{2t} .
\end{equation}
Before differentiating again, note that every factor on the right equals one at
$t=0$: $M_k(0)=1$ for all $k$, and $\ee^0=1$. That is what makes the second
derivative tractable. Differentiating the product gives a sum of $n$ terms, each
identical to \cref{m:eq:dmgf2} except that a single factor has been differentiated,
and at $t=0$ all the undifferentiated factors become one, leaving
\begin{equation*}
M_n''(0) = (2p)^n\lrbs{M_{n-1}'(0) + M_{n-2}'(0) + \cdots + M_1'(0) + 2}.
\end{equation*}
But $M_k'(0)$ is the mean population at generation $k$, namely $(2p)^k$, so
\begin{equation*}
M_n''(0) = (2p)^n\lrbs{(2p)^{n-1} + \cdots + (2p) + 2}
= (2p)^n\lrbs{\sum_{i=0}^{n-1}(2p)^i + 1}.
\end{equation*}
Summing the finite geometric series and rearranging,
\begin{equation}
\label{m:eq:secondMoment}
\ex{Z_n^2} = (2p)^n\lrb{1+\frac{1}{1-2p}} - (2p)^{2n}\lrb{\frac{1}{1-2p}} .
\end{equation}
The minus sign is not optional: at $n=1$ the expression must return
$\ex{Z_1^2}=4p$, and it does.

\subsubsection{Late-time variance}

From $\var{Z_n} = \ex{Z_n^2}-\ex{Z_n}^2$ and \cref{m:eq:secondMoment},
\begin{equation*}
\var{Z_n} = (2p)^n\lrb{1+\frac{1}{1-2p}} - (2p)^{2n}\lrb{\frac{1}{1-2p}} - (2p)^{2n},
\end{equation*}
and for $2p<1$ the terms in $(2p)^{2n}$ are negligible beside $(2p)^n$ at large
$n$, so $\var{Z_n}\to(2p)^n\bigl(1+1/(1-2p)\bigr)$. It is tempting to read
stationarity of the conditional variance straight off that, but the step is not
available: the conditional variance is not the unconditional one divided by
anything simple. Start instead from the conditional moments, using
\cref{m:eq:lateS} in each:
\begin{align*}
\ex{Z_n^2\mid Z_n>0} &= \frac{\ex{Z_n^2}}{S_n}
\;\xrightarrow{n\to\infty}\;
\frac{(2p)^n\lrb{1+\frac{1}{1-2p}}}{A(p)(2p)^n}
= \frac{1}{A(p)}\lrb{1+\frac{1}{1-2p}},\\[4pt]
\ex{Z_n\mid Z_n>0} &= \frac{\ex{Z_n}}{S_n}
\;\xrightarrow{n\to\infty}\;
\frac{1}{A(p)} .
\end{align*}
Both limits are finite, so the conditional variance tends to a constant:
\begin{equation}
\label{m:eq:condvar}
\var{Z_n\mid Z_n>0} \;\longrightarrow\;
\frac{1}{A(p)}\lrb{1+\frac{1}{1-2p}} - \frac{1}{A(p)^2} .
\end{equation}
The first two moments of the conditional population therefore both stabilise, and
the same argument carried to higher moments gives the further moments of the
quasi-stationary law. Note that \cref{m:eq:condvar}, like \cref{m:eq:qsmean}, is
expressed entirely in terms of $p$ and the single unknown constant $A(p)$.
Everything now turns on that constant.

\subsection{What is known about \texorpdfstring{$A(p)$}{A(p)}, and where}
\label{m:sec:Apinterface}

Two facts about $A(p)$ are used freely in the remainder of this chapter, and both
are established in \ChA:
\begin{enumerate}[label=(\roman*)]
\item the limit \cref{m:eq:lateS} exists and satisfies $0<A(p)\le\tfrac12$ on
      $[0,\tfrac12)$, so that the quasi-stationary mean $1/A(p)$ is at least $2$;
\item $A(p)$ is computable to arbitrary accuracy by iterating \cref{m:eq:SnIter},
      which is how the numerical values quoted in \cref{m:fig:condmean} and below
      were obtained.
\end{enumerate}
Beyond that, this chapter treats $A(p)$ as a named object and nothing more. Its
exact infinite-product and series representations, the two-sided bounds that trap
it between the continuous-time constant $\Ac(p)$ of \cref{m:sec:cts} and twice
that, the near-critical asymptotic $A(p)\sim2(1-2p)$ with its logarithmic
correction, the failed search for an elementary closed form, and the identification
of $A(p)$ with a value of the Koenigs linearising function of the logistic map ---
which is why that search failed, and should have been expected to --- are all the
subject of \ChA. Nothing in the present chapter depends on them.

\subsection{A continuum view of the survival recursion}
\label{m:sec:riccati}

Before leaving the discrete process it is useful to record a continuum
approximation that recurs below. Writing
\begin{equation}
\label{m:eq:epsdef}
\varepsilon = 1-2p
\end{equation}
for the distance from criticality --- a symbol used in this sense throughout the
thesis --- the recursion \cref{m:eq:SnIter} rearranges to
\begin{equation}
\label{m:eq:Sdiff}
S_{n+1}-S_n = -\varepsilon S_n - pS_n^2 .
\end{equation}
For large $n$ the unit increment is small on the scale over which $S$ varies, and
the left-hand side may be replaced by a derivative. As in
\cref{m:rem:continuumcritical} this replacement is heuristic; what it produces is
the Riccati equation
\begin{equation}
\label{m:eq:riccati}
\dda{S}{n} = -\varepsilon S - pS^2 .
\end{equation}
The substitution $u=1/S$ linearises it to $u' = \varepsilon u + p$, with solution
$u(n) = K\ee^{\varepsilon n}-p/\varepsilon$ and hence
\begin{equation}
\label{m:eq:riccatisol}
S(n) = \frac{\varepsilon}{K\varepsilon\,\ee^{\varepsilon n}-p},
\end{equation}
where $K$ is fixed by the initial condition. Setting $\varepsilon=0$ recovers the
critical scaling $S\sim2/n$ of \cref{m:eq:twoovern}; keeping $\varepsilon>0$
recovers the geometric decay $S\sim A(p)(2p)^n$ of \cref{m:eq:lateS}, since
$(2p)^n = \ee^{n\ln(1-\varepsilon)}$ agrees with $\ee^{-\varepsilon n}$ to leading
order in $\varepsilon$. The constant $K$ is exactly what the continuum
approximation cannot supply, because it is determined by the early generations,
where the approximation is invalid. That is another way of saying that $K$ is
$A(p)$, and that no argument of this kind will produce it --- an observation \ChA{}
takes as its starting point. The same Riccati structure reappears in
\cref{m:sec:absbirthdeath}, where it is the characteristic equation of a partial
differential equation and can be solved outright.

\FloatBarrier
